%%%%%%%%%%  scan idx 299  |  folio 286  %%%%%%%%%%
\textbf{2.10. Accouplement entre \(K^{\bullet}\) et \(K_{\bullet}\).} Soit \((X,\mathcal O_X)\) un topos commutativement annelé. On sait (I 5.8.4 b)) que l'on a un accouplement
\[
 \otimes_A^{\mathrm L}:D(X)_{\uparf}\times D^b(X)_{\coh}
 \longrightarrow D^b(X)_{\coh}.
\]
Celui-ci définit un homomorphisme
\begin{equation}\tag{2.10.1}
 K^{\bullet}(X)\otimes K_{\bullet}(X)\longrightarrow K_{\bullet}(X).
\end{equation}
Grâce,à l'associativité du produit tensoriel, on voit que (2.10.1) \emph{munit} \(K_{\bullet}(X)\) \emph{d'une structure de} \(K^{\bullet}(X)\)-\emph{module}. Plus généralement, si \(f:(X,\mathcal O_X)\longrightarrow(Y,\mathcal O_Y)\) est un morphisme de topos commutativement annelés, on peut regarder \(K_{\bullet}(X)\) comme un \(K^{\bullet}(Y)\)-module via l'homomorphisme \(f^*:K^{\bullet}(Y)\longrightarrow K^{\bullet}(X)\).

\textbf{2.11. Loi covariante sur \(K_{\bullet}\) et formule de projection.} Soit \((X,A)\) un topos annelé. À la différence de \(K^{\bullet}({}_A X)\), le groupe \(K_\bullet({}_A X)\) n'a pas de variance naturelle par rapport à \((X,A)\). Cependant, en Géométrie Algébrique et en Géométrie Analytique, \(K_{\bullet}(X)\) a tendance à être \emph{covariant} en \(X\) pour les morphismes \emph{propres}. Précisons ce point :

\textbf{2.11.1.} Soit \(f:X\to Y\) un morphisme propre et pseudo-cohérent (III 1.2) de schémas quasi-compacts. Alors, sous réserve de validité de la conjecture de finitude pour \(f\) (III 2.1), donc en particulier (III 2.2) si \(Y\) est noethérien (auquel cas l'hypothèse de pseudo-cohérence sur \(f\) est automatiquement vérifiée (III 1.3)) ou si \(f\) est projectif, le foncteur \(Rf_*\) induit, en vertu de (III 2.6), un foncteur
\[
 Rf_*:D^b(X)_{\coh}\longrightarrow D^b(Y)_{\coh}.
\]
Celui-ci définit un
homomorphisme
\begin{equation}\tag{2.11.1.1}
 f_*:K_{\bullet}(X)\longrightarrow K_{\bullet}(Y)
\end{equation}
